\begin{table}[htbp]
\centering
\caption{Robustness and Heterogeneity}
\label{tab:robustness}
\begin{adjustbox}{max width=\textwidth}
\begin{tabular}{lccccc}
\toprule
 & (1) & (2) & (3) & (4) & (5) \\
 & Rural & Urban & High Poverty & Low Poverty & Dose-Response \\
\midrule
CS Exposure $\times$ Post & -0.0363 & 0.0017 & -0.1045 & 0.0104 & \\
 & (0.0965) & (0.0712) & (0.1158) & (0.0640) & \\
Medium CS $\times$ Post & & & & & -0.1991* \\
 & & & & & (0.1059) \\
High CS $\times$ Post & & & & & -0.1339 \\
 & & & & & (0.1327) \\
\midrule
County FE & Yes & Yes & Yes & Yes & Yes \\
State $\times$ Year FE & Yes & Yes & Yes & Yes & Yes \\
\midrule
RI $p$-value & \multicolumn{5}{c}{0.112} \\
LOSO range & \multicolumn{5}{c}{[-0.0900, -0.0238]} \\
\midrule
Observations & 15,800 & 15,794 & 14,898 & 16,475 & 31,620 \\
\bottomrule
\end{tabular}
\end{adjustbox}
\begin{tablenotes}[flushleft]
\small
\item \textit{Notes:} Columns (1)--(2) split the sample by median county population (2017). Columns (3)--(4) split by median poverty rate (2017). Column (5) replaces the continuous CS Exposure with tercile indicators (Low CS share is the omitted category). RI $p$-value is from 500 permutations of the treatment assignment. LOSO shows the range of the main coefficient when each state is excluded. All specifications include county and state $\times$ year fixed effects. Standard errors clustered at the state level. * $p<0.10$, ** $p<0.05$, *** $p<0.01$.
\end{tablenotes}
\end{table}

